\section{Methods}
\label{sec:methods}

The analytics component of HospitAI consists of an offline pipeline with three stages plus the AI component. Stage one consumes the exports and generates a clean relational database (ETL process). Stage two computes the patient features and clusters patients via the K-Prototypes algorithm. Stage three computes the clinical sequences in each patient group. Above all this, the AI component will take in the preprocessed information and generate readable text output, either as automatically generated cluster/sequence descriptions or as a question-answer session for humans. We explain the details below.

\subsection{Stage 1: Data Ingestion and ETL}

The data flow in HospitAI begins with the comma-separated exports generated by Synthea and ends with a clean PostgreSQL database.

Three key decisions related to reduction and data cleansing influence the results of our analysis. First, we fix one hospital. In particular, the organization which was selected has the highest utilization. Only busy providers are kept. Three hundred providers are ranked by decreasing number of encounters. Then all encounters of the excluded providers are redistributed between other providers with a fixed random seed so that the number of encounters and the provider load remain unchanged. Second, we fix the dates. Some events created by Synthea may have the timestamp set to some date in the future. Therefore, we restrict the timestamp to the current date. We obtain the patient's age by subtracting the birth date from the death date (or from the current date for patients who are still alive). The length of an encounter is estimated using the beginning and end of the encounter time and limited by some reasonable value depending on the encounter's class (for instance, up to two hours for the wellness visit and thirty days for the inpatient encounter). Third, missing values are not mean-imputed, since that would distort the numerical centroids K-Prototypes treats as real measurements. Sparse categorical and vital-sign features instead fall into an explicit ``Unknown'' level, and count features are filled with zero, which correctly represents an absence of events.

The number of records in the observation table is in millions, so the ingestion is done in parts. This step produces normalized data that includes patients, encounters, conditions, medications, procedures, observations, care plans and providers.

\subsection{Stage 2: Feature Engineering and Clustering}
For clustering, we need a table in which every row represents a patient and every column represents a feature. Most effort here goes into creating such a table, since the clinically relevant information is scattered across the database tables and not directly available for modeling. For each feature, a separate SQL aggregation query is run.

We start off with a large collection of potential features and end up selecting a few non-overlapping ones later. These features are grouped according to type in Table~\ref{tab:features}. The numeric features are relatively simple (counts, lengths and ratios). On the other hand, there is more domain knowledge embedded within the categorical features because we do not rely on raw codes, which are derived into clinically meaningful tiers through rule-based SQL. Disease features such as diabetes, hypertension, heart disease, cancer and kidney disease are defined based on keyword matching against condition descriptions, specifically excluding ``prediabetes'', ``suspected'' or ``family history'' when the feature would imply otherwise. Care behavior features include a tier of thirty-day readmissions generated using window functions on consecutive visits, a flag for those who make visits to the ED in out-of-office hours, their drug use pattern which integrates their diagnostic information with the results of the AUDIT-C and DAST-10 tests, and a tier for different pain management approaches.

\begin{table}[t]
\caption{Engineered patient features, grouped by kind. The full catalogue is computed offline. Clustering uses a non-redundant subset.}
\label{tab:features}
\centering
\footnotesize
\begin{tabular}{@{}p{0.26\linewidth}p{0.66\linewidth}@{}}
\toprule
Group & Features \\
\midrule
Numerical & age, active conditions, total medications, total procedures, distinct providers seen, lifetime medication days, encounter duration hours, preventive-to-reactive ratio \\[2pt]
Demographic & gender, marital status, race, ethnicity \\[2pt]
Disease flags & diabetes, hypertension, heart disease, cancer, kidney disease, mental health, respiratory profile \\[2pt]
Binned vitals & BMI class, blood-pressure risk, glucose risk (each with an explicit ``Unknown'' level) \\[2pt]
Care-pattern tiers & dominant care plan (severity-ranked), condition burden, medication intensity, polypharmacy, 30-day readmission tier, out-of-hours flyer \\[2pt]
Risk profiles & substance-use profile, pain-management tier, cholesterol risk, pregnancy history, smoking status \\
\bottomrule
\end{tabular}
\end{table}

The selected engineering methods include the following. Some features are heavily skewed, so we apply a natural log transform and then min-max scaling. For the physiological vital sign features such as BP, Glucose, we assume that they may not be available, hence categorizing them along with unknowns instead of replacing the missing values with mean values as part of retaining the knowledge that the particular feature was never recorded. The most critical care plan is not the most common care plan but the care plan that is most severe. We arrange the care plans per patients according to specific severity rankings (oncology and palliative, maternal, trauma, respiratory and cardio-metabolic), then choose the most severe care plan. So a cancer patient is assigned the cancer care plan even though their most common care plan can be a checkup. We do not transform categorical features through one hot encoding because K-Prototypes can work with them directly. There is also the clinical severity score (first principal component of log-transformed counts of procedures/medications/conditions).

K-Prototypes algorithm \cite{huang1998} is an extension of K-Means algorithm for mixed data. Distance is defined by the sum of squares of the difference of numerical features and the number of mismatches in categorical features. The weighted sum using weight $\gamma$ helps to control the relative significance of the categorical features. We use the Huang initialisation and select the value of $\gamma$ automatically using the library.

Prior to the model fitting, we look for any statistical outliers based on a z-score criterion for the numerical variables after log-transforming the heavily skewed counts, such that the tail distribution itself does not lead to the detection of an outlier. If a z-score greater than six is detected for any of the variables, the corresponding patient is considered an outlier and there are 289 such patients that constitute around 1\% of the dataset. None of the outliers are removed from the analysis. Instead, they receive a unique label of $-1$ and are included in the analysis with the other patients without being used in the K-Prototypes cluster fitting, i.e., none of the clusters are generated for them. This is done because their values would skew the cluster generation, but they are clinically important since they represent the most severe cases like cancer, dialysis, and heart disease.

\subsection{Stage 3: Clinical Sequence Mining}

In the third phase, we study the order of clinical events. One ordering sequence is defined for each patient, i.e., the ordering sequence of all the events occurring in a certain timeline, arranged according to the dates. Each event is prefixed with a type character such that the source table for the particular event is known (Condition is labeled C, Medication M, Procedure P, and Care plan CP). If more than one event occurs on a particular day, then the ordering is based on a clinical order (C-M-P-CP).

Cleaning comes before mining, where we remove all events except for those occurring during the last two years per patient, exclude some purely administrative events and psychosocial ``status'' events added to almost any encounter by Synthea, merge immediate duplicates of events, and restrict the length of the event timeline, allowing us to perform mining very efficiently. Mining of sequences is done using the PrefixSpan algorithm \cite{pei2001}, a classical approach that builds on finding closed frequent sequences by projecting the data onto all the already discovered prefixes. We ask PrefixSpan to generate closed sequences of length ranging from three to six events having a support of at least one percent of patients, and we mine more than needed for selection purposes.

\begin{figure}[t]
\centering
\includegraphics[width=0.8\linewidth]{figs/sequence_mining_pipeline.png}
\caption{The sequence-mining pipeline}
\label{fig:pipeline}
\end{figure}

The key step happens around PrefixSpan in Figure~\ref{fig:pipeline}. When ranking the patterns by their supports only, the best patterns are going to be very frequent but not useful sequences, such as regular dentistry cleaning actions which everyone undergoes. To solve this, we rank the patterns using Geometric Mean Lift (GML). Given a pattern $p$ containing $k$ elements with its support being $s(p)$ among $n$ patients and item supports being $s(i)$, we compute
\begin{equation}
\mathrm{GML}(p) \;=\; \frac{s(p)/n}{\left(\prod_{i \in p} s(i)/n \right)^{1/k}} .
\end{equation}
It is just another step in applying the notion of lift from two items to $k$ items based on the geometric mean calculation. It can be calculated in log scale to prevent overflow. The higher GML value is, the more probable it is that the pattern would appear in practice more often than its component frequencies' product.

Two other concepts have been proposed. Firstly, an additional bonus is assigned to the overall score as multiplication by a multiplier of $(1 + 0.3\, (t-1))$ where $t$ stands for the total number of different types of events in the pattern. In this way, we give preference to the actual trajectory of a patient which includes diseases, drugs, and surgeries to a technical pattern of three drug applications. Secondly, certain diversity criteria are used. We filter out patterns including triplets of the same type of events, patterns with very similar items according to the Jaccard distance, and limit the number of patterns including the same beginning or ending events or dominated by dentistry events. Also, some slots are reserved for fully representative patterns that include all four types of events.

Finally, all these processes are run for each scope once. Mining is performed both on the global scope and on the individual clusters, including the outlier group denoted $-1$. Patterns discovered in the global scope are ignored while mining the individual clusters so as to keep only those patterns that define the clusters. Further, each of the patterns retained has their patient ID attached to them using the intersecting posting list of the least frequent items.

\subsection{AI Interpretation Layer}

Finally, this layer makes the output understandable. We use just one LLM using the API of OpenAI. However, we do this wisely: the statistics are never computed by the machine alone. The programming code computes everything before providing these static values to the machine in its prompt, asking it to verbalize about no other data but the provided numbers. Such an approach is known as ``verbalizer'' since it guarantees that facts are described in the text.

The key point is that clinical features are given as auxiliary data within the prompt. Clusters' prompts are constructed out of actual figures: population averages, cluster-specific values with an annotation ``compared to population average'', prevalence rates of various conditions, averages for vital indicators, dominant care plan category, and prevalent diseases, therapies, and interventions. In turn, the model is supposed to generate strict JSON including the risk level, clinical reasoning for doctors, statement for patients, key characteristics of a certain cohort or sequence, and suggestions as to what actions should be taken. Sequences interpretation is done analogically by providing mined patterns accompanied by aggregated patient profile (age, gender, BMI, disease prevalence, dominant care plan), and further decoding the prefixes and temporal arrows into reasonable care pathway.

All AI outputs are cached as a JSON value together with their corresponding rows, so the model is run only once per row.

\subsection{Natural-Language-to-SQL Agent}

In addition to strict interpretations, HospitAI allows users to ask questions on top of the available data. A question-answer agent was developed, which is a state machine based on LangGraph, shown in Figure~\ref{fig:agent}. Each state accomplishes one task alone.

\begin{figure}[t]
\centering
\includegraphics[width=1\linewidth]{figs/agent_diagram.png}
\caption{The natural-language-to-SQL agent}
\label{fig:agent}
\end{figure}

In the initial phase the question is classified and rephrased. The agent determines whether the question requires the database, whether it is a generic medical question that can be answered directly, or whether it is irrelevant and must be rejected.

In case the database needs to be accessed, the query generation procedure is performed. Generation is based on the embedding of the database schema and a set of rules. Generated queries are validated in two layers. The first, deterministic layer ensures the query is a single read-only statement and blocks potentially harmful keywords via regex validation. The second layer validates the query semantics, i.e. its ability to answer the question with the appropriate tables and filters. If the query fails either layer, the error message is passed back to the agent to regenerate the query until an attempt limit is reached. Database execution errors are treated the same way.

The agent is also role-aware. If the request comes from a doctor, an additional scope condition restricts every search to the records of that doctor's own patients via their provider id, and any request without this condition is rejected, preventing a doctor from seeing other doctors' patients.